\section{Introduction}
\label{sec:introduction}
% Motivation, problem statement, goals, research question

Chess has long been regarded as a canonical benchmark problem in artificial
intelligence and algorithm engineering. Despite its simple and fully specified
rules, the game exhibits an enormous search space caused by its high branching
factor and strategic depth. As a consequence, exhaustive search of the complete
game tree is computationally infeasible, even with modern hardware, and chess
engines must rely on heuristic search techniques to make strong decisions under
limited computational resources.

At the core of most classical chess engines lies a depth-limited game tree
search based on Minimax or Negamax, enhanced by Alpha-Beta pruning. Over the
years, a variety of improvements have been proposed to increase search
efficiency, including iterative deepening, transposition tables, quiescence
search, and heuristic move ordering. Among these techniques, iterative
deepening has become a standard approach, particularly in time-controlled
settings, as it allows the engine to gradually refine its decision while always
maintaining a valid best move.

The primary motivation of this work is to investigate the effectiveness of
iterative deepening when operating under explicit time constraints. While
iterative deepening is widely used in practice, it repeatedly re-searches the
game tree at increasing depths, which raises the question of how much efficiency
is lost compared to a fixed-depth search when the available time is limited.
This leads to the central research question of this thesis:

\emph{To what extent does the effectiveness of iterative deepening degrade under
strict time constraints, and how well can modern search heuristics compensate
for this overhead?}

To address this question, a complete chess engine is implemented from scratch in
\texttt{C++}. The engine deliberately employs a classical mailbox-based board
representation rather than highly optimized bitboard techniques. This design
choice emphasizes architectural clarity, ease of reasoning, and suitability for
academic analysis, allowing the impact of search algorithms and time management
to be studied in isolation from low-level optimizations.

The implemented engine integrates a Negamax search with Alpha-Beta pruning,
iterative deepening, quiescence search, transposition tables, and heuristic move
ordering. A dedicated time management module supports multiple search modes,
including fixed depth, fixed time, node-limited, and infinite search, following
the UCI protocol. In particular, a two-stage time budget with soft and hard
deadlines is used to control the interaction between iterative deepening and
time constraints.

The contribution of this work is twofold. First, it provides a clean and modular
implementation of a classical chess engine that serves as a transparent
experimental platform. Second, it presents an empirical evaluation of
iterative deepening under time constraints, analyzing its impact on search
depth, node counts, and move stability. The results aim to shed light on the
practical trade-offs between search efficiency and decision quality in
time-controlled chess engines.

The remainder of this report is structured as follows. \cref{sec:related_work}
discusses relevant prior work in computer chess and game tree search.
\cref{sec:preliminaries} introduces the necessary theoretical background.
The architecture and implementation of the engine are described in
\cref{sec:architecture} and \cref{sec:search}. The evaluation function and time
management are detailed in subsequent sections, followed by an experimental
evaluation and a discussion of the results. Finally, conclusions and directions
for future work are presented.